\chapter{The Cantor Set}
\label{appendix:cantor_set}

The Cantor ternary set is a classical counterexample in analysis and topology, illustrating a set that is uncountable, compact, nowhere dense, and has measure zero.

\section{Construction of the Cantor Set}
Start with the closed unit interval $C_0 = [0,1]$.
\begin{enumerate}
    \item Remove the open middle third $(1/3, 2/3)$, leaving $C_1 = [0, 1/3] \cup [2/3, 1]$.
    \item Remove the open middle third of each of the two remaining intervals, leaving $C_2 = [0, 1/9] \cup [2/9, 1/3] \cup [2/3, 7/9] \cup [8/9, 1]$.
    \item Repeat this process indefinitely. At step $n$, $C_n$ consists of $2^n$ closed intervals, each of length $3^{-n}$.
\end{enumerate}
The \textbf{Cantor set} $C$ is the intersection of all $C_n$:
\[
C = \bigcap_{n=0}^\infty C_n.
\]

% TikZ Diagram of Cantor Set Iteration
\begin{figure}[htbp]
\centering
\begin{tikzpicture}[scale=10, thick]
    % C0
    \draw[mainblue, ultra thick] (0, 0.8) -- (1, 0.8) node[right] {$C_0$};
    
    % C1
    \draw[mainblue, ultra thick] (0, 0.6) -- (1/3, 0.6);
    \draw[mainblue, ultra thick] (2/3, 0.6) -- (1, 0.6) node[right] {$C_1$};
    
    % C2
    \draw[mainblue, ultra thick] (0, 0.4) -- (1/9, 0.4);
    \draw[mainblue, ultra thick] (2/9, 0.4) -- (1/3, 0.4);
    \draw[mainblue, ultra thick] (2/3, 0.4) -- (7/9, 0.4);
    \draw[mainblue, ultra thick] (8/9, 0.4) -- (1, 0.4) node[right] {$C_2$};
    
    % C3
    \foreach \x in {0, 2/9, 2/3, 8/9} {
        \draw[mainblue, ultra thick] (\x, 0.2) -- (\x + 1/27, 0.2);
        \draw[mainblue, ultra thick] (\x + 2/27, 0.2) -- (\x + 1/9, 0.2);
    }
    \node[mainblue, right] at (1, 0.2) {$C_3$};
    
    % C4
    \foreach \x in {0, 2/27, 2/9, 8/27, 2/3, 20/27, 8/9, 26/27} {
        \draw[mainblue, ultra thick] (\x, 0) -- (\x + 1/81, 0);
        \draw[mainblue, ultra thick] (\x + 2/81, 0) -- (\x + 1/27, 0);
    }
    \node[mainblue, right] at (1, 0) {$C_4$};

    % Ticks on C0
    \draw (0, 0.78) -- (0, 0.82) node[above] {$0$};
    \draw (1, 0.78) -- (1, 0.82) node[above] {$1$};
    \draw (1/3, 0.58) -- (1/3, 0.62) node[above] {$\frac{1}{3}$};
    \draw (2/3, 0.58) -- (2/3, 0.62) node[above] {$\frac{2}{3}$};
\end{tikzpicture}
\caption{Step-by-step construction of the Cantor ternary set $C = \bigcap_{n=0}^\infty C_n$.}
\label{fig:cantor_set}
\end{figure}

\section{Properties of the Cantor Set}
\begin{theorem}
The Cantor set $C$ satisfies:
\begin{enumerate}[label=\rm{(\roman*)}]
    \item $C$ is compact (closed and bounded in $\R$).
    \item $C$ is nowhere dense in $\R$ ($\Int(C) = \emptyset$).
    \item $C$ is perfect (it has no isolated points, $C' = C$).
    \item $C$ is uncountable ($\card(C) = c = 2^{\aleph_0}$, via ternary expansion $x = \sum_{n=1}^\infty \frac{a_n}{3^n}, a_n \in \{0,2\}$).
    \item $C$ has Lebesgue measure zero ($\lambda(C) = 1 - \sum_{n=0}^\infty \frac{2^n}{3^{n+1}} = 0$).
\end{enumerate}
\end{theorem}


\chapter{Limits in Metric Spaces}
\label{appendix:limits}

Let $(X, d_X)$ and $(Y, d_Y)$ be metric spaces, $A \subseteq X$, $f: A \to Y$, and $x_0 \in A'$.

\begin{definition}[Limit of a Function]
We say that $L \in Y$ is the \textbf{limit of $f(x)$ as $x$ approaches $x_0$}, written $\lim_{x \to x_0} f(x) = L$, if for every $\eps > 0$, there exists $\delta > 0$ such that:
\[
0 < d_X(x, x_0) < \delta \implies d_Y(f(x), L) < \eps.
\]
\end{definition}

\begin{theorem}[Sequential Characterisation of Limits]
$\lim_{x \to x_0} f(x) = L$ if and only if for every sequence $(x_n)$ in $A \setminus \{x_0\}$ with $x_n \to x_0$, $f(x_n) \to L$.
\end{theorem}
